\section{近期研究（2026 年 8--11 月）}
\paragraph{Tri-DehazeGS：分离雾介质，并补偿被低透射率压弱的训练信号}
2026.09~\cite{2609_11223}。Tri-DehazeGS 将场景辐亮度与跨视图共享的三维雾介质分开表示。它同时补偿介质透射率降低后被压弱的重建梯度，使低可见区域能向场景表示提供更充分的约束。

论文在 RealX3D、Mip-NeRF 360 和 Fog-NeRF 比较，并在 RealX3D 上相对所列最强散射感知基线 3D-UIR 报告 PSNR 高 3.35 dB。重建仍依赖估计深度正则，细结构或低纹理处的深度过度平滑会导致边界模糊。

\paragraph{RGS：在反射重建中约束跨视图几何并按反射增密}
2026.09~\cite{2609_19421}。镜面亮斑会随视角移动，直接用球谐颜色拟合容易以不一致的表面几何解释变化。RGS 基于二维高斯做延迟着色，用 VGGT 的多视图几何先验约束形状一致性，再按反射线索增密高斯以表达高光。几何正则与着色承担不同任务。

论文在 Shiny Blender、Glossy Synthetic、Ref-NeRF Real 和 NeRF Synthetic 共 23 个场景报告 PSNR、SSIM 和 LPIPS。该指标主要回答反射物体的新视图质量，不能用来断言材料或几何已正确分解。对几何先验误差较大的捕获场景，还需检验这种正则的适用范围。

\paragraph{GS-PI：以条件点云扩散分解高斯资产材质}
2026.09~\cite{2609_19907}。GS-PI 先以已有重建的几何为基础，用受多视图信息约束的点云扩散预测材质，再执行短时蒸馏，以较保守的方式调整不透明度和尺度。这将逐场景的光照与材质联合拟合转化为学习式分解，便于从已有高斯生成可重光照资产。

主文的材质和重光照比较使用真值前景掩码，重光照还共享真值法线；补充材料另列各系统使用自身法线的端到端比较。输入重建缺失时，点云条件也缺失，复杂镜面反射与细小结构仍可出现误差。

\paragraph{从多视图中分离并重建镜头眩光}
2026.09~\cite{2609_39527}。镜头眩光是随相机和光源改变的成像效应，将它写进静态场景高斯会把伪影误当作世界中的物体。作者先微调潜空间扩散模型得到去眩光视图，再用多视图共同拟合场景与眩光。眩光用位于场景前方二维平面上可移动的高斯表示，渲染时与场景层一起合成。

论文在 Flare7K++ 和自建 VFX 基准上评价二维去眩光，并另做多视图眩光重建和消融。二维去除指标与成像分解的质量要分别判断。镜头效应能否由前置平面模型表示，也取决于具体光学结构与观测条件。

\paragraph{EvenSplat：同时解释图像曝光与场景照明}
2026.10~\cite{2610_01876}。单视图全局曝光补偿无法消除空间不均匀照明，例如同一幅图里的亮区和阴影区。EvenSplat 以图像空间分解提供逐视图约束，并在高斯层共享三维照明成分。全局相机响应和局部补偿用于解释剩余的逐视图误差，减少它们对场景外观的干扰。

作者在实拍及受控合成的相机曝光变化、空间照明变化和高对比度照明设置中评估，优势主要出现在不均匀及高对比度条件。原文也列出 sRGB 的饱和和分解歧义、局部补偿依赖训练图像，以及未建模的镜面反射、运动模糊和噪声。

\paragraph{VolS-GS：在高斯的空间支撑内模拟次表面光传播}
2026.10~\cite{2610_04007}。局部材质模型难以描述光线进入半透明介质后在别处出射的现象。VolS-GS 在高斯覆盖的体积内学习散射和吸收，结合表面着色与次表面项合成图像，并从同一传输过程提取阴影线索。这样在分解高光、遮挡和内部散射时，共享了对场景空间结构的约束。

论文在 SSS-GS、NRHints 和 GS3 的单灯照明基准上比较重光照，提升最稳定的是次表面散射占主导的数据。该结果依据其体积传输近似和单灯采集协议，不能直接推断任意照明与透明介质下的准确性。

\paragraph{Casual Flash：结合闪光与环境光观测分解材质}
2026.10~\cite{2610_06035}。静态环境光下，材料颜色和未知照明会产生歧义。该工作把随相机移动的闪光与不闪光的室内图像作为互补观测，在高斯逆渲染里分别建模固定和闪光照明，并使用跨视图一致的漫反射表示来约束材质分解。

论文用非配对闪光与不闪光的合成和真实室内数据比较材质估计和重光照。作者指出，相对逐灯采集协议，这种简化采集对细微粗糙度变化的敏感度较低，过曝光区域的观测也不能提供同样的约束。

\paragraph{PAM-ToD：冻结街景几何，按时段适配乘性光照与加性发光}
2026.10~\cite{2610_11572}。从白天重建切到夜间时，部分变化近似整体曝光或色温，部分变化则来自新出现的灯光。PAM-ToD 保留原街景高斯的权重，利用少量有位姿的目标时段锚点图像拟合乘性颜色校正与非负加性发光项。与材质及照明完全分解相比，该插件只处理源时段到目标时段的外观变化。

作者以 StreetGS 为底座，在 CARLA-ToD 静态和动态条件使用一组或八组同步锚点比较，并提供 Waymo 定性例子。重建和新视图轨迹使用不同源图采样规则，不能混用结果。仅改变颜色的插件也不能替代目标时段新出现物体的重新建模。

\paragraph{以可见场景几何约束近场反射}
2026.10~\cite{2610_13757}。近场物体的反射随表面位置变化，仅用固定远场环境图容易把可见的空间结构抹平。该方法以二维高斯表达表面，将用于直接渲染的主高斯也引入反射线查询，让直接观测到的结构约束反射。混合外观分量与视角相关高光模型处理不能由直接观测充分解释的外观。

原文在 NeRF-Casting、EnvGS 和 Ref-Real 的已有数据及自建 NearRef 数据上检验新视图效果，以图像指标和法向可视化分别报告不同层次的结果。法向示例不是覆盖全部场景的定量几何验证。作者还将更高效的光线跟踪列为后续工作。

\paragraph{VisibilitySplat：把稀疏射线可见性引入高斯逆渲染}
2026.10~\cite{2610_19847}。光源被遮挡时，光照变化可能被错误吸收到表面材质中。VisibilitySplat 使用 BVH 加速的点基射线查询估计透射率，用少量查询结果补充常规高斯光栅化训练，为材质与光照分离提供显式遮挡信号。

论文在合成基准评价新视图、反照率与重光照。消融中显式可见性的收益主要出现在分解和重光照，纯 PBR 变体仍可能获得更高的新视角外观分数。不同指标反映的目标应分开比较。

\paragraph{DaFISH：按相机运动旋转球谐照明}
2026.10~\cite{2610_24811}。当光源随相机移动时，常规高斯会把逐帧照明变化固化在场景颜色中。DaFISH 将发射颜色分为反照率和方向性照明，并用 Wigner-D 旋转把球谐系数随相机位姿变换。场景仍由高斯表征，变化发生在随视图变化的光照项。

论文在 DarkRobotic 四个场景及 FourLights 序列进行对照，报告五个场景的 SSIM 与 LPIPS 优势。FourLights 缺少 DarkGS 所需的光源标定，因而该方法未进入这项对照。结果对应相机光源共同运动的采集条件。

\paragraph{GraftRay：用估计的面光源把网格插入高斯场景}
2026.10~\cite{2610_37036}。GraftRay 给高斯赋予估计的材质和法线，从用户标记的窗户或灯具初始化面光源，再通过高斯感知的路径追踪与单位光照基底拟合光源辐亮度。差分合成将插入网格引发的直接光、接触阴影与反弹光变化叠加到原场景。

论文在室内采集和合成场景比较插入效果，并讨论整图指标对局部照明变化不敏感的问题。方法需要用户标记光源区域，二维材质和法线先验的错误也会传给光源拟合和物体光照。

\paragraph{NIRJALA：为水下高斯加入前向散射和可逆约束}
2026.11~\cite{2611_23450}。NIRJALA 从高斯渲染辐亮度与逆深度，分别预测水中的直达、后向和前向散射三条光路。训练时以低维物理模型对水下图像做正反向变换约束柔性的神经水体模型，水体分解和去水后的颜色平衡由此分别受约束。

论文的评价区分真实水下图像重现与有参考的去水、几何恢复。SeaThru-NeRF 的真实捕获没有对应的干净图像或真值深度，对这些场景的水下新视角分数不能替代清晰颜色或几何的验证。

\paragraph{用分离的辐射高斯场渲染多波段星系}
2026.11~\cite{2611_24723}。该方法为恒星辐射、尘埃消光与尘埃辐射分别拟合高斯场，以物理基函数约束各场的光谱，并通过分阶段优化重建多波段视图。高斯表示承载各物理分量，不把同一像素的所有能量当成一种颜色系数。

评价使用 SKIRT-Auriga 和 TNG50 的模拟数据。其中 Auriga 子集为十个星系，每个以十七视图中的十四个训练、三个测试。作者指出，尘埃吸收的星光和其红外辐射未通过模型显式耦合，因此物理一致性仍有边界。

\paragraph{MergeableRF：直接合并外观网络以合成新光照}
2026.11~\cite{2611_25523}。MergeableRF 对不同已知照明下的外观网络使用共享初始化和合并感知训练，让多个光源模型的参数可组合成单一外观网络。这样可直接用合并后的场为 NeRF 或 3DGS 渲染新照明，避免对多个光照各渲染一次再在图像上相加。

实验使用 Neural Gaffer 合成物体与 OLATverse 的真实逐灯捕获。方法依赖已知的训练灯光和可组合的光照基底，不能直接套用到只有一种未知光照的普通视频采集。
